\chapter{Sackgassen: was nicht funktioniert hat und warum}
\label{ch:sackgassen}

% Quellen: verworfen.md X1-X9, counterexamples.md, lab-notes.
% Ton: erzählend, ehrlich, pro Sackgasse eine klare Lehre.

Forschungsberichte erzählen gern nur die Treffer. Dieses Kapitel
macht das Gegenteil: Es sammelt die Wege, die wir gegangen und
wieder zurückgegangen sind. Nicht aus Bußfertigkeit, sondern weil
die Sackgassen dieses Projekts mehr über die Horimetrik verraten
haben als mancher Satz -- und weil mindestens eine von ihnen sich
später als Umweg statt Irrweg entpuppte.

\section{Das $n!/2$-Debakel, oder: drei Datenpunkte sind kein Muster}

\begin{sackgasse}
\emph{Idee:} Die Schalenspringer-Folge zeigte bei $n=5,6,7$ die
Werte $60, 360, 2520$ -- exakt $n!/2$. Eine Formel!
\emph{Hindernis:} Der vierte Datenpunkt: $21168\ne20160$.
\emph{Lehre:} Muster aus wenigen Termen sind Hypothesen, keine
Ergebnisse.
\end{sackgasse}

Diese Geschichte hat ein bemerkenswertes Nachspiel: Viel später
fiel eine exakte Ziffernformel für dieselbe Folge vom Himmel -- und
sie erklärt rückwirkend, \emph{warum} die falsche Formel drei Terme
lang stimmte. Für alle drei Treffer reduziert sich die
Ziffernformel auf einen einzigen Summanden, und der beträgt jeweils
genau $n!/2$: bei $120$ und $5040$, weil sie in ihrer eigenen
Schale Ein-Ziffern-Darstellungen haben, bei $720$, weil die erste
Maximalziffer die Formel schon nach dem ersten beitragenden
Summanden abschneidet. Die Koinzidenz war kein Rauschen,
sondern ein echtes Strukturmerkmal -- nur eben eines der Ziffern
von $n!$, nicht der Folge selbst. So rehabilitieren sich Sackgassen
manchmal: Der Fehler von damals wurde zur Fußnote eines Satzes.

\section{Die Kleine-Welt-Falle, sechsmal}

Dieselbe Falle hat uns im Lauf des Projekts sechs Mal erwischt, in
immer neuen Verkleidungen: Eine Schranke des Kompaktifizierungs-Kommutators
sah bis zur siebten Welt wie $\pm1$ aus (sie ist in Wahrheit
beidseitig unbeschränkt); ihre Nachfolgerin, eine einseitige
Schranke, hielt Zufallsstichproben bis zur einundzwanzigsten Welt
stand und fiel erst durch \emph{konstruierte} Zeugen; ein Element
des Kompaktifizierungsmonoids wirkte bis zur vierten Welt
idempotent, und der kleinste Gegenzeuge wohnt in der fünften; und
das $n!/2$-Muster kennen Sie schon. Die Fälle fünf und sechs
ereigneten sich bei der oberen Seite des Tiefengesetzes an einem
einzigen Abend: Erst starb eine hübsche Dreieckszahlen-Hypothese am
eigenen Kandidaten-Sweep, noch ehe sie aufgeschrieben war -- dann
entpuppte sich die Schwelle, die dieser Sweep lieferte, ihrerseits
als Artefakt einer zu flachen Suche. Die gemeinsame Lehre ist
präziser als das übliche „mehr testen“:

\begin{center}
\emph{Im Horiraum wachsen die interessanten Gegenbeispiele
fakultätisch langsam.\\ Wer Schranken behauptet, muss konstruieren,
nicht stichproben.}
\end{center}

Übrigens spukt durch mehrere dieser Geschichten dieselbe Zahl:
die Sechs -- Ziel des kleinsten Schalensprungs (die Fünf, um eine
Null verlängert, wird zur Sechs), wertkleinster Trennzeuge des
Monoids, Anfang der ersten interessanten Schale. Sie ist der
Hausgeist dieses Buchs.

\section{Verworfene Anwendungsträume}

\begin{sackgasse}
\emph{Idee:} Hori-Adressen statt IPv6, Hori-Heaps statt Binärheaps.
\emph{Hindernis:} Der Horiraum ist innerhalb eines Horizonts
strukturell starr -- Adressräume brauchen lokale Elastizität, und
der Heap zahlt für die wachsende Verzweigung mit teuren
Kindervergleichen.
\emph{Lehre:} Die Stärke des Systems ist die globale
Fakultätsordnung, nicht lokale Flexibilität. Anwendungen müssen zur
Struktur passen, nicht umgekehrt.
\end{sackgasse}

\begin{sackgasse}
\emph{Idee:} Mächtigkeits-Argumente -- ist der Horiraum „größer“
als die natürlichen Zahlen?
\emph{Hindernis:} Er ist abzählbar; jede Nummerierung plättet ihn.
\emph{Lehre:} Der Mehrwert liegt nicht in der Kardinalität, sondern
in Struktur pro Wert. Ein Schachbrett wird nicht interessanter,
wenn man seine Felder durchnummeriert -- und nicht langweiliger.
\end{sackgasse}

Der dritte Anwendungstraum kam von außen, war der ehrgeizigste --
und starb am ersten Satz dieses Buchs.

\begin{sackgasse}
\emph{Idee:} Kryptographie aus Kapazität (externer Vorschlag,
Juli 2026). Eine geheime Gestalt wird durch $k$ führende Nullen
geschoben und wertkompakt ausgewertet; veröffentlicht wird nur der
nackte Wert. Der Angreifer, so die Hoffnung, scheitert am
Differenzüberlauf -- der Rückweg über Anti-Differenzen sprengt die
Kapazität -- und der Seitentreue-Satz versperrt jede hybride
Abkürzung.
\emph{Hindernis:} Der Rückweg braucht weder Anti-Differenzen noch
gemischte Arithmetik. Aus $y=P(N)$ holt fortgesetzte Division mit
Rest die Ziffern einzeln herunter: $c_0=y\bmod N$ ist \emph{exakt}
(denn $c_0\le r<N$), dann liefert $y_1=(y-c_0)/N$ die nächste Runde
mit Divisor $N{-}1$, und so fort -- eine Handvoll Divisionen, und
die „geheime“ Gestalt liegt offen. Das ist kein raffinierter
Angriff, sondern die Bijektion aus Kapitel~\ref{ch:horizonte} am
verschobenen Auswertungspunkt: Die Gewichte $1,N,N(N{-}1),\ldots$
bilden wieder ein Mischbasensystem, und die Ziffernschranken
bleiben brav unter den Divisoren.
\emph{Lehre:} Die Fakultätsordnung verwirbelt nicht, sie
\emph{sortiert} -- dieselbe Bijektivität, die das System so sauber
macht, verbietet ihm die kryptographische Härte. Zwei
Präzisierungen fürs Protokoll: Der Seitentreue-Satz verbietet
gemischte \emph{Halbringe}, nicht gemischte \emph{Algorithmen} --
ein Angreifer darf dividieren, tabellieren und Repräsentationen
wechseln, ohne je eine Arithmetik zu mischen. Und eine
Einwegfunktion braucht ein Geheimnis; die Horimetrik hat keins, sie
ist mit Absicht überall gleich hell. Was bleibt, ist die
bescheidenere Rolle: Das Horizontkalkül aus
Kapitel~\ref{ch:strukturpolynome} taugt als \emph{Messinstrument}
für Kapazitäts- und Rauschwachstum in fremden Konstruktionen --
nicht als eigene Härtequelle.
\end{sackgasse}

\section{Die tropische Fata Morgana -- und ihre Auflösung}

Die lehrreichste Sackgasse verdient den Schluss. Früh fiel auf,
dass die Horizonte sich unter Addition wie ein Maximum und unter
Multiplikation wie eine Summe verhalten \emph{wollen} -- das
Verhalten der sogenannten tropischen Algebra. Der direkte Versuch,
so zu rechnen, scheiterte an einem Überlaufproblem: Die Wertsumme
kann den Maximum-Horizont sprengen. Verworfen.

Monate später -- in der Zeitrechnung dieses Projekts: Stunden --
kehrte die Idee durch die Hintertür zurück. Der Horizont zerfällt
in einen Pflichtanteil und einen frei wählbaren Überschuss, und
\emph{auf dem Überschuss} funktioniert die tropische Arithmetik
tadellos: Der Pflichtanteil fängt alle Überläufe, das Gedächtnis
rechnet mit Maximum und Summe. Die Idee war nie falsch. Sie stand
nur im falschen Koordinatensystem.

\begin{oberflaeche}
Wenn dieses Kapitel eine Botschaft hat, dann diese: Verwerfen ist
keine Niederlage, sondern Buchführung. Von unseren neun
dokumentierten Sackgassen (das Arbeitsregister führt sie als X1
bis X9) wurde eine später zum Satz -- die tropische --, eine zur
Fußnote eines Satzes -- das $n!/2$-Debakel --, und die übrigen
wurden zu präzisen Warnschildern.
Nichts davon wäre passiert, hätten wir die Fehlversuche gelöscht
statt datiert.
\end{oberflaeche}
